\Qhead{3.4}{Отклик линейной системы на шумовые воздействия}{Отклик линейной системы на шумовые воздействия; функция Грина, интеграл Дюамеля. Действие шума на колебательный контур, фильтрация шума.}
\Subfile{3.4.1}{sections/03_statistics/q04/s01}

\recap{Временная свёртка и умножение спектров — два описания одной линейной фильтрации. Дисперсия выхода равна площади входного спектра, взвешенного квадратом передачи.}
\board{Причинная свёртка, $S_y=|H|^2S_x$, резонансный спектр и равновесная площадь.}
\oral{Почему при уменьшении трения тепловой резонансный пик растёт, а средняя энергия остаётся той же?}
\bait{Можно оставить шумовую полосу RC-фильтра: она определяется площадью $|H|^2$, а не только точкой половинной мощности.}
\sources{Биленко и соавт.~\cite{bilenko}, §§2.1.5--2.1.6, 4.1.7, 4.2.1; Митрофанов~\cite{mitrofanov}, лекции о тепловых флуктуациях. Равновесные дисперсии согласованы с ФДТ в вопросе~\Qref{3.6}.}
